La question touche au statut explicatif des constantes. Lorsqu'une
théorie reçoit un paramètre déterminé expérimentalement, ses équations
établissent les conséquences découlant de cette valeur. Une dérivation
structurelle modifie cette relation : la valeur devient l'une des
conséquences de la construction. La question de l'intensité d'un
couplage ou d'un rapport caractéristique de la matière devient
alors celle des relations mathématiques qui les déterminent.
Le passage du paramètre à la conséquence constitue l'argument directeur
de cet article.

La détermination numérique acquiert, dans cette perspective, un contenu
constitutif. L'équation réunit des opérations dont la provenance doit être
explicitée ; son évaluation fournit la valeur que ces opérations fixent. Une
fois établies les conditions correspondantes d'existence, d'unicité et de compatibilité,
la nécessité du résultat relève de la démonstration.
Expliquer pourquoi une constante prend une valeur signifie reconstruire cette
nécessité dans le domaine de la théorie. La précision décimale manifeste
une conséquence quantitative de la structure, tandis que la dérivation
expose la raison pour laquelle cette conséquence précise est obtenue.

La portée de cette question s'accroît lorsque plusieurs déterminations
partagent une même origine. La thèse proposée par les auteurs relie
les constantes à l'organisation produisant leurs relations réciproques.
L'architecture commune doit expliquer tant l'individualité de chaque
coordonnée que les dépendances qui la rattachent aux autres.
L'ordre des constructions acquiert ainsi une valeur démonstrative :
il permet de reconnaître quelles données sont engendrées, quelles opérations les utilisent et
quelles relations subsistent dans les réalisations successives. La diversité
des valeurs est étudiée comme expression d'une unité formelle articulée.

Dans cette organisation, la métrologie remplit une fonction ultérieure et
essentielle. La construction fixe ses résultats avant de les confronter aux
grandeurs observées ; l'expérience examine leur correspondance
physique et délimite la portée de l'accord empirique. La nécessité mathématique s'établit
relativement aux prémisses et aux opérations déclarées, et la réalisation
naturelle est examinée par l'identification des observables et leur
confrontation indépendante. Les deux tâches contribuent à une explication
complète : la première détermine le résultat au sein du corpus ; la seconde
étudie sa pertinence pour décrire la constitution quantitative du
monde.

Cette perspective permet de situer historiquement la recherche. Les
antécédents considérés ci-dessous précisent différentes manières de
relier lois, paramètres et observation. Le présent travail adopte
comme problème central la détermination interne des valeurs et de la
structure qui les relie. Ses développements arithmétiques, géométriques et
algébriques tirent leur unité argumentative de cette question. L'aspiration
ontologique consiste à expliquer comment des propriétés physiques apparemment
indépendantes peuvent correspondre à des relations nécessaires d'une même
organisation mathématique.
